\chapter{Controlling Execution}\label{ch:gdb-exec}

A debugger is a program that runs your program slowly enough for you to look at
it. This chapter is the controls: start, stop, and advance by one step of
whatever size you need.

\section{Starting and Stopping}\label{sec:gdb-run}

\begin{keys}[0.28]
	\cmd{run}                 & Start the program from the beginning                 \\
	\cmd{run arg1 arg2}       & \dots{} with arguments                               \\
	\cmd{run < in.txt > out.txt} & \dots{} with redirection, as in \autoref{sec:sh-redirect} \\
	\cmd{start}               & \cmd{run}, with a temporary breakpoint on \cmd{main} \\
	\cmd{continue}            & Resume until the next stop                           \\
	\cmd{continue 5}          & Resume, ignoring this breakpoint the next 4 times     \\
	\cmd{kill}                & Stop the program, keep the session                   \\
	\cmd{quit}                & Leave GDB                                            \\
	\cmd{set args -v in.txt}  & Set the arguments for future runs                    \\
	\cmd{show args}           & What they currently are                              \\
\end{keys}

\begin{note}
	\cmd{start} is almost always the right way in: it gets you to the first line of
	\cmd{main} with the program loaded and the dynamic libraries mapped, which is
	the point at which most things become inspectable. Plain \cmd{run} with no
	breakpoints set simply runs to completion.
\end{note}

\begin{note}[The program keeps its state]
	\cmd{run} a second time restarts from scratch; breakpoints and watchpoints
	survive, convenience variables survive, and the value history survives. So the
	loop is: set breakpoints, \cmd{run}, look, adjust, \cmd{run} again --- without
	ever leaving GDB.
\end{note}

\section{Stepping}\label{sec:gdb-step}

\begin{keys}[0.28]
	\cmd{next} (\cmd{n})      & Run the next source line. \textbf{Step over} calls   \\
	\cmd{step} (\cmd{s})      & Run the next source line. \textbf{Step into} calls   \\
	\cmd{finish}              & Run until the current function returns; print its value \\
	\cmd{until}               & Like \cmd{next}, but do not go backwards --- runs a loop to completion \\
	\cmd{until 60}            & Run until line 60 is reached, or the frame returns   \\
	\cmd{advance parse}       & Run until \cmd{parse} is reached, or the frame returns \\
	\cmd{next 10}             & Ten lines at once                                    \\
\end{keys}

\begin{intuition}
	\cmd{next} and \cmd{step} differ only at a function call. \cmd{next} treats the
	call as one line and runs the whole thing; \cmd{step} stops at its first line.
	Use \cmd{next} by default and \cmd{step} only when you have decided that the
	bug is inside \emph{this} call --- otherwise you will spend the afternoon
	inside \cmd{printf}.
\end{intuition}

\begin{eg}
\begin{gdbcode}
(gdb) start
Temporary breakpoint 1, main () at main.c:12
12          int n = count_words(argv[1]);
(gdb) next                     # run count_words, stop after it
13          printf("%d\n", n);
(gdb) print n
$1 = 0                         # wrong: the file has 40 words
(gdb) run                      # start again, this time go inside
(gdb) step
count_words (path=0x7fffffffe5c1 "input.txt") at parser.c:8
8           FILE *f = fopen(path, "r");
\end{gdbcode}
	The first \cmd{next} establishes \emph{that} the answer is wrong; the
	\cmd{step} begins finding out why. Doing it in that order saves you from
	stepping through functions that were never the problem.
\end{eg}

\begin{note}[\cmd{step} into a library]
	\cmd{step} will not enter a function with no debug information --- it steps
	over it instead, which is what you want. If you have installed libc's debug
	symbols and now find yourself inside \cmd{malloc}, \cmd{finish} gets you out,
	and \cmd{skip file malloc.c} stops it happening again.
\end{note}

\begin{eg}[\cmd{until}, and why loops are annoying]
	Stepping with \cmd{next} inside a loop walks the body once per iteration, so a
	thousand-iteration loop takes a thousand presses. \cmd{until} at the closing
	line of the body runs the loop out:
\begin{gdbcode}
(gdb) next
14              for (int i = 0; i < n; i++) {
(gdb) next
15                  total += weight[i];
(gdb) until                    # run the loop to completion
17          printf("%d\n", total);
\end{gdbcode}
	The other answer is a conditional breakpoint on the iteration you care about
	--- \autoref{sec:gdb-cond}.
\end{eg}

\begin{moral}
	\key{<CR>} repeats the last command, so stepping is: type \cmd{n} once, then
	hold Enter. This is the single most-used key sequence in GDB.
\end{moral}

\section{Source Lines versus Instructions}\label{sec:gdb-instr}

Everything above steps by \emph{source line}. Underneath, a line is several
machine instructions, and sometimes that is the level you need.

\begin{keys}[0.28]
	\cmd{stepi} (\cmd{si})  & One machine instruction, entering calls        \\
	\cmd{nexti} (\cmd{ni})  & One instruction, stepping over calls           \\
	\cmd{disassemble}       & The current function's instructions            \\
	\cmd{disassemble /s}    & \dots{} interleaved with the source            \\
	\cmd{x/5i \$pc}         & The next five instructions from the program counter \\
	\cmd{info registers}    & Every register                                 \\
	\cmd{display/i \$pc}    & Show the next instruction after every stop     \\
\end{keys}

\begin{eg}
\begin{gdbcode}
(gdb) disassemble /s main
Dump of assembler code for function main:
main.c:
5       int main(void) {
   0x1149 <+0>:     push   %rbp
   0x114a <+1>:     mov    %rsp,%rbp
6           int *p = NULL;
   0x114d <+4>:     movq   $0x0,-0x8(%rbp)
7           return *p;
   0x1155 <+12>:    mov    -0x8(%rbp),%rax
   0x1159 <+16>:    mov    (%rax),%eax        # <-- the fault is here
\end{gdbcode}
\end{eg}

\begin{note}
	You need this level in exactly three situations: the code was compiled without
	\cmd{-g}; the bug is in inline assembly or a compiler intrinsic; or you suspect
	the compiler, which is almost never the case. Otherwise stay at the source
	level --- assembly is a way of being precise about the wrong thing.
\end{note}

\section{Changing the Flow}\label{sec:gdb-force}

GDB can do more than watch. These commands alter a running program, which is
occasionally exactly what an experiment needs.

\begin{keys}[0.28]
	\cmd{return}          & Abandon the current function immediately            \\
	\cmd{return 0}        & \dots{} making it return that value                 \\
	\cmd{jump 42}         & Continue from line 42 instead of here               \\
	\cmd{call parse(s)}   & Call a function in the program, right now           \\
	\cmd{print parse(s)}  & The same, and show the result                       \\
\end{keys}

\begin{eg}
	Testing a fix without recompiling. A function is returning a null pointer and
	you suspect the caller mishandles it:
\begin{gdbcode}
(gdb) break lookup
(gdb) run
Breakpoint 1, lookup (key=0x4006f4 "port") at config.c:31
(gdb) return 0                 # force the NULL return, right now
(gdb) continue
Program received signal SIGSEGV, Segmentation fault.
0x00005555555552a1 in main () at main.c:22
22          printf("%s\n", value);
\end{gdbcode}
	Confirmed in twenty seconds: the caller does not check. You have reproduced a
	rare failure on demand, which is usually the hard part.
\end{eg}

\begin{note}
	\cmd{call} and \cmd{print} of a function actually run it in the program's own
	process, side effects and all. That is powerful --- \cmd{print dump\_state()}
	invokes your own debug printer --- and dangerous: if the function crashes, it
	crashes the program you were debugging. GDB will tell you and offer to unwind.
\end{note}

\begin{moral}
	\cmd{jump} skips the stack adjustment a real control transfer would do, so
	jumping out of one function into another is undefined nonsense. Jump within a
	function, or not at all.
\end{moral}
